
où, si \(M\) est un
espace vectoriel sur \(k\), on écrit \(M_X\) pour \(M\otimes_k{O}_X\). Pour \(n\ge m\), \(\alpha_n\) est surjectif par (ii); donc, si \(\beta_{n+1}\) est surjectif, alors \(\beta_n\) l'est. Par le théorème de Serre, \(\beta_{n+1}\) est surjective pour \(n\gg0\); d'où (iii).

\textbf{Proposition 1.4.} \emph{Soit \(0\to F(-1)\to F\to G\to0\) une suite exacte sur \(X\). Si \(G\) est \(m\)-régulier, alors on a}~:
\begin{enumerate}[label=(\roman*)]
\item \emph{\(H^q(F(n))=0\) pour \(q\ge2\) et \(n\ge m-q\).}
\item \emph{\(h^1(F(n-1))\ge h^1(F(n))\) pour \(n\ge m-1\).}
\item \emph{\(H^1(F(n))=0\) pour \(n\ge (m-1)+h^1(F(m-1))\).}
\end{enumerate}
\emph{En particulier, \(F\) est \([m+h^1(F(m-1))]\)-régulier.}

En effet, dans la suite exacte
\[
H^{q-1}(G(n))\longrightarrow H^q(F(n-1))\longrightarrow H^q(F(n))\longrightarrow H^q(G(n)),
\]
les deux groupes extrêmes sont \(0\) pour \(q\ge2\) et \(n\ge m-(q-1)\) par 1.3 (i). Donc \(H^q(F(m-q))\xrightarrow{\sim}H^q(F(m-q+1))\xrightarrow{\sim}H^q(F(m-q+2))\xrightarrow{\sim}\ldots\), mais, par le théorème de Serre, \(H^q(F(n))=0\) pour \(n\gg0\); d'où (i).

Pour \(n\ge m-1\), considérons la suite exacte suivante :
\[
0\to H^0(F(n-1))\to H^0(F(n))
  \xrightarrow{\alpha_n}H^0(G(n))\to H^1(F(n-1))
  \to H^1(F(n))\to0
\]
Elle implique (ii). Or, supposons \(\alpha_n\) surjectif et considérons le diagramme commutatif suivant :
\[
\begin{tikzcd}[column sep=large,row sep=large]
H^0(F(n))\otimes H^0({O}_X(1))
  \arrow[r,"\alpha_n\otimes\mathrm{id}"]
  \arrow[d]
& H^0(G(n))\otimes H^0({O}_X(1))
  \arrow[r]
  \arrow[d,"\beta_n"]
& 0 \\
H^0(F(n+1))
  \arrow[r,"\alpha_{n+1}"]
& H^0(G(n+1)).
\end{tikzcd}
\]
